% !TEX root = main.tex
\subsubsection{Independent biases and scalar perturbations}
\label{app:CR-independent}
For mobility bounds $0<\nu_-\le\nu_+\le1$, define
\begin{equation}
 \begin{aligned}
 0<\initscale\le\widehat\initscale(\nu_-,\nu_+):=\min\Biggl\{&\frac12,
 \frac{\sqrt{\nu_-/\nu_+}}{2(m+2)},
 \left(\frac{\nu_-}{16(\log m+2)}\right)^{1/4},\\
 &\frac1{8e^{x_\star}\sqrt{d_{\mathrm h}}(m+2)(\log(4x_\star\sqrt{d_{\mathrm h}})+4)}\Biggr\}.
 \end{aligned}
 \label{eq:CR-variable-scale}
\end{equation}
\begin{corollary}[Independent Euclidean distractor biases]
\label{cor:CR-independent}
Replace each distractor score $-B$ by $-B_j$, train every $B_j$ independently by ordinary Euclidean gradient flow,
and initialize $B_j(0)\in[0,1]$. Keep the other hypotheses of part (ii).
Define
\[
 \mathcal D=\sum_{j=1}^m e^{-B_j},\qquad B_{\mathrm{eff}}=\log(m/\mathcal D).
\]
For $0<\initscale\le\widehat\initscale(1/m,\min\{1,e/m\})$, the same time bound \cref{eq:CR-time} gives
entry into $x\ge x_\star$, $\mathcal D\le1$, retention until $L\le z$, and persistent positive margins on every binary transition.
In particular the biases need not start equal or remain equal, and the $\initscale^{-4}$ polynomial rate is retained.
\end{corollary}
\begin{proof}
The complete variable-mobility and independent-bias argument is retained in the
\href{archival_full_version.pdf}{accompanying archival supplement}, Appendix D.3.1.
The loss depends on the biases through $\mathcal D=\sum_j e^{-B_j}$; the reciprocal
identity $\ddt{e^{B_j}}=\pd{L}{\mathcal D}$ gives mobility in
$[1/m,\min\{1,e/m\}]$ and the sufficient scale stated there.
\end{proof}
\paragraph{Gaussian scalar perturbations.}\label{cor:CR-scalar-Gaussian}
The Gaussian scalar/offset specialization and its quantitative tail conditions
are retained with their proof in the same archival section.
